% ATLAS: starting visual testing -- the seven jobs for the first time anyone opens the Phase 1 code.
% Written by hand; compile with pdflatex. The look (fonts, colours, grey command blocks, tick boxes, note boxes) is the
% one of the Station PC guide. Facts are the team lead's, as of 20:25 on 3 October 2026.
\documentclass[10pt]{article}
\usepackage[a4paper,top=13mm,bottom=17mm,left=16mm,right=16mm,footskip=9mm]{geometry}
\usepackage[T1]{fontenc}
\usepackage[utf8]{inputenc}
\input{glyphtounicode}\pdfgentounicode=1
\usepackage{helvet}
\renewcommand{\familydefault}{\sfdefault}
\usepackage[varqu,varl]{zi4}
\usepackage[table]{xcolor}
\usepackage{array,booktabs}
\usepackage{enumitem}
\usepackage{needspace}
\usepackage{listings}

\pdfinfo{/Title (ATLAS: starting visual testing)}

% ---------- colours (the palette of the other ATLAS handouts) ----------
\definecolor{ink}{HTML}{1F2933}
\definecolor{mute}{HTML}{56616B}
\definecolor{laneD}{HTML}{E3EDFC}
\definecolor{accC}{HTML}{9A4A00}
\definecolor{accD}{HTML}{1B55B0}
\definecolor{cmdbg}{HTML}{F1F3F5}
\definecolor{trapbg}{HTML}{FDF3E7}
\definecolor{laterbg}{HTML}{EEF2F6}

% ---------- footer ----------
\makeatletter
\def\ps@atlas{%
  \def\@oddhead{}\def\@evenhead{}%
  \def\@oddfoot{\parbox{\textwidth}{\footnotesize\color{mute}ATLAS: starting visual testing\hfill 3 October 2026\hspace{6mm}Page \thepage}}%
  \let\@evenfoot\@oddfoot}
\makeatother
\pagestyle{atlas}

\setlength{\parindent}{0pt}
\setlength{\parskip}{0pt}
\color{ink}

% ---------- commands to type, and what they print ----------
% Commands are set so that copying them out of the PDF gives exactly what is printed:
% plain ASCII quotes and hyphens, one command per line, no prompt characters, no wrapped lines.
\lstdefinestyle{cmd}{basicstyle=\ttfamily\fontsize{8.35}{11.3}\selectfont\color{ink},columns=fullflexible,
  keepspaces=true,upquote=true,showstringspaces=false,backgroundcolor=\color{cmdbg},frame=none,
  aboveskip=3.5pt,belowskip=3.5pt,xleftmargin=6pt,framexleftmargin=6pt,breaklines=false}
\lstdefinestyle{inl}{basicstyle=\ttfamily\fontsize{8.8}{11}\selectfont,columns=fullflexible,keepspaces=true,
  upquote=true,showstringspaces=false,breaklines=true}
\lstnewenvironment{cmd}{\lstset{style=cmd}}{}
\newcommand{\cc}{\lstinline[style=inl]}
\newcommand{\code}[1]{{\ttfamily\fontsize{8.8}{11}\selectfont #1}}
\newcommand{\see}{\par\vspace{0.6mm}{\footnotesize\color{mute}You should see:}}

% ---------- jobs and their steps ----------
\newlength{\cbw}      % width of a tick box and its gap; measured after \begin{document}
\newcommand{\jobletter}{}
% \Job{letter}{title}{one line: why}   the blue band of the Station guide's parts, with the letter of the job
\newcommand{\Job}[3]{%
  \par\vspace{6mm}\needspace{52mm}\gdef\jobletter{#1}%
  {\setlength{\fboxsep}{0pt}\colorbox{laneD}{\parbox{\linewidth}{\vspace{2.2mm}\hspace{3mm}%
    {\fontsize{17}{19}\selectfont\bfseries\color{accD}#1}\hspace{3mm}%
    {\fontsize{13.5}{16}\selectfont\bfseries\color{ink}#2}\vspace{2.2mm}}}}\par\nopagebreak
  \vspace{1.8mm}{\small\color{mute}#3\par}\nopagebreak\vspace{1.4mm}}
% the steps of a job: a tick box, then the job letter and a number (A1, A2, ...)
\newlist{jsteps}{enumerate}{1}
\setlist[jsteps]{label=\cb\textbf{\jobletter\arabic*},leftmargin=12.5mm,labelwidth=10.4mm,labelsep=2.1mm,align=left,itemsep=2.8mm,topsep=1.5mm,parsep=0.8mm}
\newcommand{\act}[1]{\textbf{#1}\enskip}
% a step keeps its first lines and its first command on one page
\newcommand{\step}{\needspace{21mm}\item}
% ticked sub-steps inside one step
\newlist{ticks}{itemize}{1}
\setlist[ticks]{label=\cb,leftmargin=\cbw,labelsep=0mm,align=left,itemsep=1.6mm,topsep=1.2mm,parsep=0.6mm}

% ---------- boxes: a trap, and a thing that waits ----------
\newsavebox{\notebox}
\newenvironment{trap}[1][If it fails]{\par\vspace{1.4mm}\begin{lrbox}{\notebox}%
  \begin{minipage}{\dimexpr\linewidth-5mm\relax}\small\raggedright{\bfseries\color{accC}#1}\enskip}
  {\end{minipage}\end{lrbox}{\setlength{\fboxsep}{2.5mm}\colorbox{trapbg}{\usebox{\notebox}}}\par\vspace{1.4mm}}
\newenvironment{later}[1][Waits for code]{\par\vspace{1.4mm}\begin{lrbox}{\notebox}%
  \begin{minipage}{\dimexpr\linewidth-5mm\relax}\small\raggedright{\bfseries\color{accD}#1}\enskip}
  {\end{minipage}\end{lrbox}{\setlength{\fboxsep}{2.5mm}\colorbox{laterbg}{\usebox{\notebox}}}\par\vspace{1.4mm}}

% ---------- small furniture ----------
\newcommand{\cb}{\raisebox{-0.25ex}{\setlength{\fboxsep}{0pt}\setlength{\fboxrule}{0.7pt}\fcolorbox{mute}{white}{\rule{0pt}{1.75ex}\rule{1.75ex}{0pt}}}\hspace{2mm}}
\newcommand{\h}[1]{\par\vspace{4mm}{\fontsize{11.5}{14}\selectfont\bfseries\color{ink}#1}\par\vspace{1.4mm}}
\newcommand{\m}{\,\textrm{\guilsinglright}\,}
% one row of the overview: a tick box, then the letter of the job
\newcommand{\jobrow}[1]{{\normalsize\cb}\textbf{\color{accD}#1}}

\begin{document}
\raggedright
\settowidth{\cbw}{\cb}

{\fontsize{27}{31}\selectfont\bfseries ATLAS: starting visual testing\par}
\vspace{2.5mm}
{\large\color{mute}The first time anyone looks at the Phase 1 code: in the Unity Editor on the Mac, then on a phone.\par}
\vspace{2mm}
{\small\color{mute}State at 20:25 on 3 October 2026. About 800 logic tests pass without the engine. Unity has not compiled
any of this code (its last compile was at 16:36, before the code existed), the six Atlas shaders have never been
compiled, and nothing has been rendered. The folder is on the branch \code{phase1-acd}.\par}

\h{The seven jobs}
%% OVERVIEW TABLE: header row "Job" and "Why", then these rows
{\small
\renewcommand{\arraystretch}{1.45}
\begin{tabular}{@{}l@{\hspace{3mm}}>{\bfseries\raggedright\arraybackslash}p{50mm}@{\hspace{3mm}}>{\raggedright\arraybackslash}p{108mm}@{}}
\toprule
 & Job & Why\\
\midrule
\jobrow{A} & Tell the Phase 1 session & It is writing files in the same folder, and it knows what you should see.\\
\jobrow{B} & Open the project in Unity & The first real compile of the scripts.\\
\jobrow{C} & Check the six shaders & A shader error shows only in Unity.\\
\jobrow{D} & Run the tests in the Editor & Five test groups need the engine and have never run.\\
\jobrow{E} & First look in play mode & The first rendered frame.\\
\jobrow{F} & Keep the evidence and report & The session can fix only what it is told.\\
\jobrow{G} & On a phone & Gate G1, first light. Needs signing first.\\
\bottomrule
\end{tabular}}

% =====================================================================
\Job{A}{Tell the Phase 1 session}{It is writing files in the same folder, and it knows what you should see.}
\begin{jsteps}
\step \act{Paste these two lines into the session that works on the branch \code{phase1-acd}:}
\begin{cmd}
I am opening this branch in the Unity Editor on the Mac now. Unity will write .meta and settings files.
Tell me what I should see in play mode in Atlas_Boot today, and which checks you want from me.
\end{cmd}

\step \act{Keep its answer open.} It is the description of a correct picture for today's code.
\end{jsteps}

% =====================================================================
\Job{B}{Open the project in Unity}{The first real compile of the scripts.}
\begin{jsteps}
\step \act{Open Unity Hub and open the project \code{ATLAS}} with Editor 6000.3.25f1. If it is not in the Hub's list:
      Add\m Add project from disk, and pick the \code{ATLAS} folder of the repository. Unity now imports every new file.

\step \act{If Unity offers Safe Mode,} it has found script errors. Either choice is fine: what matters is the Console.

\step \act{Open the Console:} Window\m General\m Console. Press Clear, and look for red lines.
      \see\ no red line. Yellow warnings are not failures.

\step \act{If there is a red line,} click the first one and copy its full text from the lower pane, then go to job F.
      Later errors usually follow from the first.
\end{jsteps}
\begin{later}[Note.]
Unity writes files of its own into the folder: a \code{.meta} file for anything new, and some settings. Leave them.
The session was told in job A.
\end{later}

% =====================================================================
\Job{C}{Check the six shaders}{A shader error shows only in Unity.}
\begin{jsteps}
\step \act{In the Project window, click each shader} and read the Inspector:

      \cc!Assets/Atlas/Rendering/Shaders/Resources/AtlasSurface.shader!\par
      \cc!Assets/Atlas/Rendering/Shaders/Resources/AtlasBackground.shader!\par
      \cc!Assets/Atlas/Rendering/Stereo/Resources/AtlasLensBlit.shader!\par
      \cc!Assets/Atlas/Rendering/Portals/Resources/AtlasPortalMask.shader!\par
      \cc!Assets/Atlas/Safety/Resources/AtlasCue.shader!\par
      \cc!Assets/Atlas/Input/Resources/AtlasReticle.shader!\par
      \see\ no list of errors in the Inspector, and no Console line that starts with \code{Shader error}.

\step \act{Copy any shader error in full} for job F. It names the file and the line.
\end{jsteps}
\begin{later}[Note.]
Four include files sit beside the first two shaders: \code{AtlasCore}, \code{AtlasWarp}, \code{AtlasGates} and
\code{AtlasPortalClip}. An error in one of them is reported against the shader that includes it. The Editor compiles
the shaders for the Mac; the phone build of job G compiles them for the iPhone.
\end{later}

% =====================================================================
\Job{D}{Run the tests in the Editor}{About 800 already pass without the engine. Five groups need it and have never run.}
\begin{jsteps}
\step \act{Open the Test Runner:} Window\m General\m Test Runner.

\step \act{EditMode tab\m Run All.}
      \see\ every test green.

\step \act{PlayMode tab\m Run All,} if that tab lists any tests.
      \see\ every test green.

\step \act{Find the five engine groups in the results:} \code{BootstrapUnityTests}, \code{AtlasActionMapTests},
      \code{HomeSpaceTests}, \code{SlicerUnityTests} and \code{LookUnityTests}. They must be green too.

\step \act{For each red test,} write down its full name and the first line of its message, for job F.
\end{jsteps}

% =====================================================================
\Job{E}{First look in play mode}{The first rendered frame.}
\begin{jsteps}
\step \act{Open the boot scene.} It is the only Atlas scene; the app builds everything else in code when play starts.

      \cc!Assets/Atlas/Scenes/Atlas_Boot.unity!

\step \act{Set the Game view to the phones' screen.} In the Game view's resolution menu, add a fixed resolution of
      2796 x 1290.

\step \act{Press Play.} Compare the picture with the session's answer from job A, and tick:
\begin{ticks}
\item Nothing is magenta. Magenta means a shader failed.
\item The picture is two eye images side by side.
\item The Console stays free of red lines while it plays, and no message repeats every frame.
\item Game view\m Stats: Batches at most 300 and Tris at most 500k. That is the plan's budget for both eyes together.
\end{ticks}

\step \act{Press Play again to stop.}
      \see\ no red line on stopping.
\end{jsteps}

% =====================================================================
\Job{F}{Keep the evidence and report}{The session can fix only what it is told.}
\begin{jsteps}
\step \act{Take a screenshot of the Game view for each check:} Cmd + Shift + 4, then Space, then click the Game window.

\step \act{Move the pictures into \code{docs/evidence/visual/}} and name them by date and check, for example
      \code{20261003\_editor\_play.png}.

\step \act{Paste into the Phase 1 session:} the first red Console line in full, any shader error, the full names of
      the red tests, and one line for each tick of job E that failed.

\step \act{After its fix, repeat jobs B to E.} Stop when the results sheet at the end is all ticks.
\end{jsteps}

% =====================================================================
\Job{G}{On a phone}{Gate G1, first light. Signing and the phones must be ready first: jobs A to C of the by-hand document.}
\begin{jsteps}
\step \act{Close the project in Unity.} The build script needs it closed.

\step \act{Build with the frame-rate display on,} in Terminal, from the repository root. The team id is needed on the
      first build only.
\begin{cmd}
ATLAS_TEAM_ID=XXXXXXXXXX ATLAS_DEV=1 tools/build_ios.sh
\end{cmd}

\step \act{Start the app with the phone in your hand} and accept the camera prompt. Then tick:
\begin{ticks}
\item The frame-rate display reads 60.
\item There are two eye images, and both move when you move the phone.
\item Nothing is magenta.
\end{ticks}

\step \act{To see and record the phone's screen on the Mac:} with the phone on its cable, open QuickTime Player\m
      File\m New Movie Recording, and choose the iPhone as the camera.
\end{jsteps}
\begin{later}[Optional: let an agent drive the Editor.]
The project includes the package MCP for Unity. With the Editor open, a Claude session started on this Mac, not in the
container, can be connected to it (Window\m MCP for Unity). It can then enter play mode, read the Console and run the
tests for you.
\end{later}

% =====================================================================
\needspace{75mm}
\h{Results sheet}
{\small
\renewcommand{\arraystretch}{1.6}
\begin{tabular}{@{}>{\raggedright\arraybackslash}p{92mm}@{\hspace{4mm}}p{10mm}@{\hspace{4mm}}p{66mm}@{}}
\toprule
\textbf{Check} & \textbf{Pass} & \textbf{Note, or the screenshot's name}\\
\midrule
The scripts compile in Unity (job B) & \cb & \\
The six shaders compile (job C) & \cb & \\
EditMode tests are green (job D) & \cb & \\
PlayMode tests are green (job D) & \cb & \\
Play mode: nothing is magenta (job E) & \cb & \\
Play mode: two eye images side by side (job E) & \cb & \\
Play mode: the Console stays clean, also on stopping (job E) & \cb & \\
Play mode: Batches at most 300, Tris at most 500k (job E) & \cb & \\
Phone: 60 frames a second (job G) & \cb & \\
Phone: both eye images follow the phone (job G) & \cb & \\
\bottomrule
\end{tabular}}

\end{document}
